\section{导读：把 CEO 当作一个 MoE 大模型来提问}

本节先建立整期访谈的阅读方法。张小珺在开场里给出一个很好的框架：把李想当作一个“CEO 大模型”，并假设它是一种 MoE（Mixture of Experts，混合专家）架构。前三个回合依次调用技术专家、战略专家和组织专家；后半程则把问题从模型、车、机器人推进到能量、亲密关系、记忆和智慧。这个框架不是玩笑，它把三小时访谈从人物访谈变成了一份关于 AI 时代公司、产品和人的系统思考。

这份笔记的目标不是复述逐字稿，而是把访谈整理成可学习的结构。第一条线是 AI 技术线：人类上下文窗口、DeepSeek 最佳实践、VLA、世界模型、Agent OS、Action 和对齐。第二条线是公司战略线：规模、用户需求、技术产品、组织能力，以及理想为什么把车、机器人和操作系统放进 AGI 时代终端的想象。第三条线是人和组织线：能量、亲密关系、共同大脑、共同心脏，以及李想对“智慧”的定义。

\begin{figure}[H]
\centering
\includegraphics[width=0.96\textwidth]{ceo-moe-interview-map.png}
\caption{CEO 大模型访谈地图：把李想当作 MoE 架构，技术、战略、组织与人共同被调用。自制概念图，依据 00:01:56--00:02:20 与 00:36:18--00:36:25 对谈内容整理。}
\end{figure}

\begin{knowledgebox}{读图：为什么 MoE 是整期的组织方式}
MoE 的直觉是“不同专家负责不同能力”。访谈用这个隐喻把李想的回答分成几个专家通道：技术专家解释 VLA 和世界模型，战略专家解释规模和平台路线，组织专家解释人才密度和方法论，最后回到一个人的能量、关系和智慧。读这期时，关键不是记住所有观点，而是看这些专家通道如何互相调用。
\end{knowledgebox}

\begin{importantbox}{本期核心命题}
AI 时代的公司不只是把模型接到产品上，而是要同时重写产品形态、组织结构、行动闭环和人的关系。车、机器人、Agent OS、世界模型和亲密关系看似分散，其实都在回答同一个问题：智能如何进入真实世界并持续产生价值。
\end{importantbox}

\begin{knowledgebox}{术语消化：本期的关键概念}
\begin{tabularx}{\textwidth}{p{0.20\textwidth}p{0.30\textwidth}X}
\toprule
术语 & 解决的问题 & 在本期中的含义 \\
\midrule
MoE & 用多个专家分担模型能力 & 既指 DeepSeek 等模型架构，也被用来隐喻 CEO 的多专家能力。 \\
VLA & 从视觉和语言走向动作 & Vision-Language-Action，是理想讨论自动驾驶和机器人的核心技术路线。 \\
World Model & 预测环境变化 & 在交通场景中用于考试、生成训练数据，并可能成为自动驾驶运营系统。 \\
Agent OS & 通用 Agent 的替代表达 & 李想认为短期不会有单一通用 Agent，但会有让专业 Agent 运行的操作系统。 \\
Action & 让 AI 真正改变世界 & 当 AI 调用工具、资产、车辆和电脑时，对齐、权限和责任都变得更重。 \\
\bottomrule
\end{tabularx}
\end{knowledgebox}

\subsection{本章小结}

EP118 的难点在于跨度极大。把它按 MoE 隐喻拆开后，主线会清楚很多：技术专家回答“模型如何进入物理世界”，战略专家回答“公司如何成为下一代平台”，组织和人回答“谁来承受这种变化”。

